% Source: https://artofproblemsolving.com/community/c3h2471814_geometric__harmonic_mean
% Stable identity: aops-topic-2471814-q01
\textbf{(a)} The two parts take $\dfrac{60}{30}=2$ hours and $\dfrac{60}{60}=1$ hour, so the average speed is $\dfrac{120}{3}=\boxed{40\text{ km per hour}}$, and indeed $\dfrac{2\cdot30\cdot60}{30+60}=40$. In general, if each half has length $d$, the total time is $\dfrac du+\dfrac dv$, so the average speed is
\[\dfrac{2d}{\dfrac du+\dfrac dv}=\dfrac{2}{\dfrac1u+\dfrac1v}=\dfrac{2uv}{u+v},\]
the harmonic mean. (It is not the arithmetic mean $45$, because more time is spent at the lower speed.)

\textbf{(b)} The growth factors are $1.8$ and $1.25$, so over two years the investment is multiplied by $1.8\cdot1.25=2.25$. A constant yearly factor $g$ gives $g^2$ over two years, so $g^2=2.25$ and $g=1.5=\sqrt{1.8\cdot1.25}$, the geometric mean of the factors. The constant growth is $\boxed{50\%\text{ per year}}$ (not the average $52.5\%$ of the two rates).

\textbf{(c)} For positive $a,b$,
\[\dfrac{a+b}{2}-\sqrt{ab}=\dfrac{\left(\sqrt a-\sqrt b\right)^2}{2}\ge0,\]
with equality only when $\sqrt a=\sqrt b$, i.e. $a=b$. Using this,
\[\dfrac{2ab}{a+b}=\dfrac{ab}{\dfrac{a+b}{2}}\le\dfrac{ab}{\sqrt{ab}}=\sqrt{ab},\]
again with equality only when $a=b$. $\blacksquare$
